\clearpage\fancyhead[L]{\small GAUSS TO BODE.}\hypertarget{native-392-lower}{}\label{native-392-lower}
\begin{center}\textbf{GAUSS TO BODE.}\par\vspace{4pt}\hrule\vspace{4pt}\fontsize{10}{12}\selectfont Astronomisches Jahrbuch für 1818. pp. 229..232. Berlin 1815.\par\vspace{4pt}\hrule\vspace{4pt}\end{center}

\begin{flushright}
Göttingen, 9 August 1815.
\end{flushright}

\par\indent I am pleased to send you all my observations of the comet. I fear the extremely unfavorable weather and the moonlight will probably prevent me from adding any new ones, although otherwise the comet could still be observed quite well for several weeks.

\SourceLeaf{393}\begin{center}\fontsize{10}{12}\selectfont\setlength{\tabcolsep}{5pt}\renewcommand{\arraystretch}{1.15}\begin{tabular}{l|c|c|c}
1815 & \shortstack{Mean time\\in Göttingen} & \shortstack{Apparent right\\ascension} & \shortstack{Apparent northern\\declination}\\\hline
March 20 & $10^{\mathrm u}\;33^{\mathrm m}\;6^{\mathrm s}$ & $54^\circ\;7'\;1''$ & $39^\circ\;7'\;47''$\\
21 & $10\;11\;37$ & $54\;33\;21$ & $39\;36\;57$\\
25 & $9\;46\;38$ & $56\;28\;50$ & $41\;38\;5$\\
30 & $9\;50\;57$ & $59\;13\;2$ & $44\;10\;27$\\
April 2 & $9\;11\;55$ & $61\;2\;27$ & $45\;39\;57$\\
June 12 & $10\;53\;31$ & $165\;3\;3$ & $52\;8\;25$\\
30 & $10\;49\;12$ & $184\;39\;24$ & $39\;22\;5$\\
July 13 & $10\;57\;20$ & $194\;32\;52$ & $29\;53\;54$\\
27 & $10\;4\;59$ & $203\;0\;56$ & $20\;37\quad:$\\
29 & $10\;13\;58$ & $204\;6\;3$ & $19\;24\;13$\\
Aug. 4 & $10\;14\;7$ & $207\;14\;8$ & $15\;53\;49$\\
\end{tabular}\end{center}

\par\indent On 13 and 27 July and 4 August, the comparison stars came from \textsc{Piazzi}’s catalogue (XII, 268; 1 Boot.; XIII, 264. 265. 284). On all other days, they came from the Histoire Céleste and the Mem. 1790. All the observations were made with the ring micrometer of the 10-foot \textsc{Herschel} telescope except the one on 12 June. On that day, the comet was compared with 219 Ursae majoris \textsc{Bode} (calculated from the Mem. 1790) using the \textsc{Frauenhofer} heliometer set up on an excellent parallactic mount. The comparison used the distance and the position angle, the latter read directly from the instrument. Excellent though the results of this method may be for brighter comets, in the present case I prefer good ring micrometer observations. The faint comet vanished completely against the star, which was more than 5 times as bright, and could only be made out when one half of the objective was mostly covered.

\Needspace{23\baselineskip}\par\indent I believe I previously sent you my parabolic elements, which I had already calculated at the end of March. They differ more from the comet’s motion in the following months than is usual in such cases. As soon as I received the observation of 12 June, I used it and two observations by \textsc{Olbers}, on 6 March and 25 April, to calculate new parabolic elements, which I give here:

\Needspace{12\baselineskip}\begin{center}\fontsize{10}{12}\selectfont Passage through perihelion, 25 April 1815.\par $11^{\mathrm u}\;41^{\mathrm m}\;19^{\mathrm s}$ Mean time in Göttingen\end{center}
\begin{center}\fontsize{9}{11}\selectfont\setlength{\tabcolsep}{2pt}\renewcommand{\arraystretch}{1.15}\begin{tabular}{lcr}
Longitude of perihelion & $.\;.\;.\;.\;.\;.\;.\;.\;.\;.\;.\;.$ & $147^\circ\;35'\;55''$\\
Longitude of the ascending node & $.\;.\;.\;.\;.\;.\;.\;.$ & $82\;43\;6$\\
Inclination of orbit & $.\;.\;.\;.\;.\;.\;.\;.\;.\;.\;.\;.\;.$ & $44\;43\;13$\\
Least distance from the Sun & $.\;.\;.\;.\;.\;.\;.\;.$ & $1.23024$\\
\end{tabular}\end{center}\begin{center}Prograde motion.\end{center}

\par\indent These elements reproduce the three observations on which they are based in full and almost exactly. Even so, they differ considerably from the intervening observations (by as much as 11 minutes in right ascension). It can often happen that three complete observations are insufficient to determine the conic section. I discussed this at greater length in the last articles of the first section of the second book of the Theoria Motus C. C. But the discrepancy between the parabolic elements and the other observations showed decisively that the orbit must differ greatly from a parabola, so I thought it would be interesting to make a preliminary calculation of the orbit without relying on a hypothesis. I used 5 observations, which the resulting ellipse reproduced apart from small discrepancies. These elements (calculated around the middle of June and printed in no. 105 of the local Gel. Anz. [p. 385 of this volume]) also \SourceLeaf{394}continued to agree well up to my latest observations. Mr. \textsc{Nicolai} has already refined them once by calculating more precisely and using more observations. He obtained only small changes. I am pleased to see that our excellent \textsc{Bessel} has also obtained results that agree very closely. I will leave the further refinement in these skillful hands. The results will be a precious legacy for our grandchildren, who will celebrate the return of this most remarkable comet.

\Needspace{14\baselineskip}\begin{center}
\emph{Elliptical elements of the comet of} 1815.
\end{center}

\begin{center}\fontsize{10}{12}\selectfont Passage through perihelion, 26 April.\par $0^{\mathrm u}\;21^{\mathrm m}\;25^{\mathrm s}$ in Göttingen.\end{center}
\begin{center}\fontsize{9}{11}\selectfont\setlength{\tabcolsep}{2pt}\renewcommand{\arraystretch}{1.15}\begin{tabular}{lcr}
Longitude of perihelion & $.\;.\;.\;.\;.\;.\;.\;.\;.\;.\;.\;.$ & $148^\circ\;58'\;48''$\\
Distance from perihelion & $.\;.\;.\;.\;.\;.\;.\;.\;.$ & $1.21349$\\
Ascending node & $.\;.\;.\;.\;.\;.\;.\;.\;.\;.\;.\;.$ & $83^\circ\;26'\;21''$\\
Inclination of orbit & $.\;.\;.\;.\;.\;.\;.\;.\;.\;.\;.\;.\;.$ & $44\;30\;43$\\
Eccentricity & $.\;.\;.\;.\;.\;.\;.\;.\;.\;.\;.\;.\;.\;.\;.$ & $0.933149$\\
\end{tabular}\end{center}

\par\indent This also gives a distance at aphelion of $35.0911$ and an orbital period of $77\tfrac12$ years. This period can be off by only a few years and is probably still a little too long.

\Needspace{9\baselineskip}\par\indent I made the following three observations of Juno, which was only 10th magnitude at this year’s opposition. They may be the only observations made anywhere on this occasion.

\begin{center}\fontsize{9}{11}\selectfont\setlength{\tabcolsep}{2pt}\renewcommand{\arraystretch}{1.15}\begin{tabular}{llrccc}
1815 & March & 1. & $10^{\mathrm u}\;22^{\mathrm m}\;0^{\mathrm s}$ & $\mathrm{AR}=197^\circ\;7'\;56''8$ & $\mathrm{Decl.}=2^\circ\;19'\;3''0\;\mathrm{S.}$\\
 & & 29. & $10\;8\;57$ & $192\;31\;51.8$ & $1\;32\;38.1\;\mathrm{N.}$\\
 & April & 8. & $9\;27\;43$ & $190\;34\;51.9$ & $2\;53\;24.2$ ---\\
\end{tabular}\end{center}

\par\indent As for the \textsc{Frauenhofer} heliometer, for which I received the exceedingly fine parallactic mount last spring, a somewhat more detailed description would be hard to understand without an exact drawing. I am convinced that if the instrument could be set up very firmly, the movements of the parallactic mount could always be relied on to within $\tfrac12$ arcminute (the verniers even read to $10''$). But my present premises do not permit this. I have to put the instrument in one doorway and then another. In the meantime I have found a way to set it up in the right position anywhere, to within 1 or at most $2'$, in very little time. When the air is favorable, for example, I can then easily find Polaris, Mercury and so on in daylight.
